\documentclass[xcolor=dvipsnames,12pt]{beamer}
\usecolortheme[named=RoyalPurple]{structure}
\useoutertheme{infolines}
\usetheme{CambridgeUS}
%\usetheme[height=14mm]{Rochester}
%\hypersetup{colorlinks=true,linkcolor=yellow}
\setbeamertemplate{items}[ball]
\setbeamertemplate{blocks}[rounded][shadow=true]
\setbeamertemplate{navigation symbols}{}
\usepackage{color}
%\usefonttheme{structuresmallcapsserif}
%\usefonttheme{professionalfonts}
%\useoutertheme{default}
%\useoutertheme{infolines}


%\documentclass{beamer}
%\usecolortheme{structure}
\usepackage{graphicx}
%\usetheme{PaloAlto}
%\setbeamertemplate{items}[ball]
%\setbeamertemplate{blocks}[rounded][shadow=true]
%\useoutertheme{umbcfootline}
\newtheorem{assume}{\rlap{A}\hskip .1in}
%\newtheorem{result}{Result}[chapter]
\newcommand{\Var}  {\mbox{Var}}
\newcommand{\V} {\mbox{V}}
\newcommand{\E}  {\mbox{E}}
\newcommand{\PR}  {\mbox{P}}
\newcommand{\diag} {\operatorname*{diag}}
\newcommand{\tr} {\mbox{tr}}
\newcommand{\bmw}[1]{{\mbox{$#1$}}}
\newcommand{\MSE}{\mbox{MSE}}
\newcommand{\MSPE}{\mbox{MSPE}}


\def\bmh{{\hat \mathbf m}}
\def\bm{{\mathbf m}}
\def\bn{{\mathbf n}}
\def\mh{{\hat m}}
\def\bOmega{{\mathbf \Omega}}
\def\be{{\mathbf e}}
\def\bs{{\mathbf s}}
\def\by{{\mathbf y}}
\def\b1{{\mathbf 1}}
\def\bx{{\mathbf x}}
\def\bw{{\mathbf w}}
\def\bzero{{\mathbf 0}}

\def\bP{{\mathbf P}}
\def\bI{{\mathbf I}}
\def\bM{{\mathbf M}}
\def\bH{{\mathbf H}}
\def\bN{{\mathbf N}}
\def\bW{{\mathbf W}}
\def\bX{{\mathbf X}}
\def\bZ{{\mathbf Z}}
\def\bT{{\mathbf T}}
\def\bV{{\mathbf V}}
\def\bC{{\mathbf C}}
\def\bG{{\mathbf G}}
\def\bY{{\mathbf Y}}
\def\bD{{\mathbf D}}
\def\bS{{\mathbf S}}
\def\bR{{\mathbf R}}

\def\bbeta{\mbox{\boldmath $\beta$}}
\def\bepsilon{\mbox{\boldmath $\varepsilon$}}

% items enclosed in square brackets are optional; explanation below
\title[]{Part 7}
%\subtitle[ttt]{}
\author[Stat 801]{Comparison of Means\\
Linear Contrasts, Orthogonal Contrasts\\
Multiple Comparison Procedures\\}
\date{}

\begin{document}

\begin{frame}[plain]
  \titlepage
\end{frame}




\begin{frame}
\frametitle{Comparison of Means}
\begin{itemize}
\item In Part 6, when the hypothesis $H_0: \mu_1 = \mu_2 = \cdots = \mu_t$ was
rejected, the inference is that at least one of the $t$
population means differs from the rest. 
\item  The next question is
which means are different from others?  
\item Is $\mu_1 \neq \mu_2$? Is $\mu_6 \neq \mu_7$?
\item Is the average ($\mu_1 + \mu_2 + \mu_3)/3$
different from ($\mu_4 + \mu_5 + \mu_6)/3$? - etc.
\item  Many times our question  will not result in a simple comparison
of whether a difference like $\mu_2 - \mu_3 = 0$ or not. 
\end{itemize}
\end{frame}



\begin{frame}
\frametitle{Comparison of Means - cont.}
\begin{itemize}
\item It may be a  more complicated question that requires a comparison like 
$\mu_1 - (\mu_2 + \mu_3)/2 = 0$ to  be made.
\item Not all questions can be formulated as comparisons. 
\item To enable us  to understand what kinds of questions  can be formulated as 
comparisons, we define a special linear function of the means.
\item A {\textcolor{magenta}{comparison}} among $t$ population
means $\mu_1,\,\mu_2, \cdots \mu_t$ can be written as the linear combination:\\[.05in]
\hspace*{1in}$\ell = a_1\mu_1 + a_2\mu_2 + \cdots + a_t\mu_t$\\[.05in]
 for given numbers $a_1, a_2, \ldots, a_t$ which satisfy
$\sum_{i=1}^t a_i = 0$.
\item Let us look at some specific examples.
\end{itemize}
\end{frame}


\begin{frame}
\frametitle{Examples:} 
Suppose $t=5$ i.e., we consider the means 
$\mu_1,\, \mu_2,\, \mu_3,\, \mu_4$, and $\mu_5$.\\[.15in]
\begin{itemize}
\item {The linear combination $\ell = \mu_2 - \mu_3$ 
has $a_i$ values\\[.2in]
 \hspace*{.7in}     $a_1 = 0,\, a_2 = 1,\, a_3 = -1,\, a_4 = 0,\, a_5 = 0.$\\[.25in]  
      Note that $\sum a_i = 0$ as required for a comparison.}\\[.25in]
\item{ The linear combination  $\ell = (\mu_1 + \mu_2)/2 - (\mu_3 + \mu_4)/2$
    has $a_i$ values\\[.18in]
\hspace*{.4in}$ a_1 = 1/2,\, a_2 = 1/2,\, a_3 = -1/2,\, a_4 = - 1/2,a_5 = 0.$\\[.25in]
Again, $\sum a_i = 0$, so it is a comparison.}
\end{itemize}
\end{frame}


\begin{frame}
\frametitle{Linear Contrasts}
\begin{itemize}
\item {A point estimate of a linear combination of population means is called a   {\textcolor{magenta}{linear contrast}}, and is given by\\[.05in]
\hspace*{1in}$\hat{\ell} = a_1 \bar{y}_{1.} + a_2 \bar{y}_{2.} + a_3\bar{y}_{3.} + \cdots + a_t \bar{y}_{t.}$\\[.05in]
 with $\sum a_i = 0$.}\\[.1in] 
\item {The estimated variance of $\hat{\ell}$ is\\[.05in]
\hspace*{1in}$\hat{V}(\hat{\ell}) = s_W^2 \sum_{i=1}^t \frac{a_i^2}{n_i}$\\[.05in]
 where $n_i$ is the number of observations taken from the $i$-th
population.} \\[.1in] 
\item {To test the hypothesis $H_0: \ell = 0$ we can use the $t$ statistic\\[.05in]
\hspace*{1in}$ t = \frac{\hat{\ell}}{\sqrt{\hat{V}(\hat{\ell})}}$\\[.05in]
 with degrees of freedom = d.f. for $s_W^2$}.
 \end{itemize}
\end{frame}


\begin{frame}
\frametitle{Orthogonal Contrasts}
\begin{itemize}
\item Two contrasts $\hat{\ell}_1 = \sum_i a_i \bar{y}_{i.}$ and
$\hat{\ell}_2 = \sum_i b_i \bar{y}_{i.}$ are {\textcolor{magenta}{orthogonal}}
whenever $\sum_i a_i b_i = 0$.  This is only defined when
$n_1 = n_2 = \cdots = n_t = n$ i.e., equal sample sizes.  
\item If all linear contrasts in a set $\hat{\ell}_1, \hat{\ell}_2, \ldots, \hat{\ell}_{t-1}$
 are pairwise orthogonal, then the set is said to be  {\textcolor{magenta}{mutually orthogonal}} set of linear contrasts.
\item Given $t$ means $\mu_1, \mu_2, \ldots, \mu_t$, and sample means
$\bar{y}_{1.}, \bar{y}_{2.}, \ldots, \bar{y}_{t.}$ (all based on
the same number $n$ of observations), the
{\textcolor{magenta}{maximum number of mutually orthogonal contrasts that exist is $(t-1)$}}. 
 \end{itemize}
\end{frame}



\begin{frame}
\frametitle{Orthogonal Contrasts - cont.}
\begin{itemize}
\item Among $t$ means there are many  $(t-1)$ sets of
contrasts that are  mutually orthogonal.
\item In a maximum mutually orthogonal set \\[.05in]
\hspace*{1in} $\hat{\ell}_1,\hat{\ell}_2, \ldots, \hat{\ell}_{t-1}$,\\[.05in]
 the linear contrasts are random variables which are statistically independent.

\item Also, the treatment sum of squares SSB is equal to the sum of the
$(\hat{\ell}_i)^2$ for any mutually orthogonal set of $(t-1)$
contrasts:
$$
  \mbox{SSB } = \sum_{i=1}^{t-1} (\hat{\ell}_i)^2.
$$
 SSB has  $(t-1)\ d.f.$ corresponding the $(t-1)$
contrasts. 
 \end{itemize}
\end{frame}


\begin{frame}
Handout - Example 9.3
\end{frame}



\begin{frame}
 Consider testing the {\textcolor{magenta}{Control vs. Agents}} comparison:
$$H_0:\,\ell = 0 \quad \mbox{vs.} \quad H_a:\, \ell \neq 0$$
where $\ell = 4\mu_1-\mu_2-\mu_3-\mu_4-\mu_5$. The 
coefficients of the corresponding contrast are therefore:
\begin{center}
\begin{tabular*}{3in}{@{\extracolsep{\fill}}lllll}
$a_1$ & $a_2$ & $a_3$ & $a_4$ & $a_5$\\
4 & -1 & -1 & -1 & -1 
\end{tabular*}
\end{center}
 Thus \\[-.3in]
\begin{eqnarray*}
\hat{\ell} & = & 4\bar{y_1}-\bar{y_2}-\bar{y_3}-\bar{y_4}-\bar{y_5}\\
 & = & 4\times1.175 -1.293 -1.328 -1.415 -1.5=-.836\\[.1in]
 V(\hat{\ell}) & = & s^2_W\cdot\sum^5_{i=1}
\frac{a^2_i}{n_i}=(0.153)\left(\frac{20}{6}\right)=.051
\end{eqnarray*}
\end{frame}
 
 
\begin{frame} 
where 
$$\sum^5_{i=1}\frac{a^2_i}{n_i}=\frac{4^2}{6}+\frac{1^2}{6}+\frac{1^2}{6}+\frac{1^2}{6}+\frac{1^2}{6}=\frac{20}{6}$$
Thus 
$$
t_c=\frac{\hat{\ell}}{\sqrt{V(\hat{\ell})}}=\frac{-.836}{\sqrt{.051}}=-3.702$$
From Table the t-table, $t_{.025,\,25}=2.06$; thus we reject $H_0$ at
$\alpha=.05$ since $|t_c|>2.06$ is in the R.R.
\end{frame}



\begin{frame}
 We also note that\\[-.3in]
$$\mbox{SSC}_1=\frac{(\hat{\ell})^2}{\sum\frac{a^2_i}{n_i}}=\frac{(-.836)^2}{20/6}=.2097$$
and therefore \\[-.3in]
$$ F_c= \frac{\mbox{SSC}_1}{s^2_W}=\frac{.2097}{.0153}=13.71$$
Since $F_{.05,\,1,\,25}=4.24$, we reject  $H_0$ at
$\alpha=.05$, the same result as above. \\[.01in]
 {\textcolor{red}{Note carefully that this sum of squares and F-test can be computed instead of the t-test.}}
 \end{frame}


\begin{frame}
Now consider testing the {\textcolor{magenta}{Biological vs. Chemical}} comparison:
$$H_0:\,\frac{\mu_2+\mu_3}{2} = \frac{\mu_4+\mu_5}{2} \quad
\mbox{vs.} \quad H_a:\, \frac{\mu_2+\mu_3}{2} \neq \frac{\mu_4+\mu_5}{2} $$
Since $H_0$ is equivalent to $.5\mu_2+.5\mu_3-.5\mu_4-.5\mu_5=0$
which is equivalent to $\mu_2+\mu_3-\mu_4-\mu_5=0$, the problem is
equivalent to testing
$$H_0:\,\ell = 0 \quad \mbox{vs.} \quad H_a:\, \ell \neq 0$$
where $\ell = \mu_2+\mu_3-\mu_4-\mu_5$. Here the contrast
coefficients are
\begin{center} 
\begin{tabular*}{3in}{@{\extracolsep{\fill}}lllll}
$a_1$ & $a_2$ & $a_3$ & $a_4$ & $a_5$\\ 
0 & 1 & 1 & -1 & -1 
\end{tabular*}
\end{center}
\end{frame}


\begin{frame}
giving $\hat{\ell}=1.293+1.328 -1.415 -1.5=-.294$ and\\ 
$$ V(\hat{\ell})=s^2_W\cdot\sum^5_{i=1}
\sum^5_{i=1}\frac{a^2_i}{n_i}=(0.153)\left(\frac{4}{6}\right)=.0102$$
where \\
$$\frac{a^2_i}{n_i}=\frac{0^2}{6}+\frac{1^2}{6}+\frac{1^2}{6}+\frac{1^2}{6}+\frac{1^2}{6}=\frac{4}{6}$$
Thus \\[-.3in]
$$
t_c=\frac{\hat{\ell}}{\sqrt{V(\hat{\ell})}}=\frac{-.294}{\sqrt{.0102}}=-2.91$$
Since $t_{.025,\,25}=2.06$; thus we reject $H_0$ at
$\alpha=.05$ since $|t_c|>2.06$ is in the R.R. 
\end{frame}


\begin{frame}
Similar to the previous comparison we may use an F-test:
$$\mbox{SSC}_2=\frac{\hat{\ell}}{\sum\frac{a^2_i}{n_i}}=\frac{(-.294)^2}{4/6}=.1297$$
and therefore 
$$ F_c= \frac{\mbox{SSC}_2}{s^2_W}=\frac{.1297}{.0153}=8.47$$ which
leads us to the same result as the t-test as  $F_{.05,\,1,\,25}=4.24$.\\[.05in]
{\textcolor{red}{The computatons for testing the other two  comparisons are similar and are not included here.}}
\end{frame}

%***************************** End new insert *****************************

\begin{frame}
\frametitle{Multiple Comparison Procedures}
\begin{itemize} 
\item We know, of course, that the
{\textcolor{magenta}{overall error rate}} when we make multiple tests is
larger than $\alpha$ (and possibly much larger). \\[.1in]
\item To compensate
for this, several different multiple comparison procedures
have been proposed to control various error rates related to
the overall error rate. \\[.1in] 
\item There are several of these; we will consider the 
following:\\[.2in]
\begin{itemize}
\item Fisher's LSD Procedure\\[.1in]
\item Tukey's W Procedure\\[.1in]
\item Scheffe's Procedure
\end{itemize}
\end{itemize}
\end{frame}


\begin{frame}
\frametitle{Multiple Comparison Procedures - cont.}
\begin{itemize} 
\item These are different procedures that each controls a different
kind of error rate and each is more or less conservative than others.  
Each has its set of fans among researchers.\\[.1in]
\item Each procedure is constructed to control a certain kind of
error rate and it is important for a user to be aware of what error
rate is controlled by a procedure before using it.  We will try to 
state how conservative each one is as we discuss it.
\end{itemize}

\noindent \textbf{\textcolor{blue}{Fisher's Protected LSD Procedure}}\\[.1in]
 The procedure is used for making all possible comparisons between
pairs of means $H_0: \mu_i - \mu_j = 0$ \ vs. \ $H_a: \mu_i - \mu_j
\neq 0$.  It presumes we rejected $H_0: \mu_1 = \mu_2 = \cdots =
\mu_t$.\\[.1in]
\end{frame}



\begin{frame}
\begin{itemize}
\item For equal sample sizes $n_1 = n_2 = \cdots = n_t = n$, consider
the $t$-test of the hypothesis above.
$$
  t = \frac{(\bar{y}_{i.} - \bar{y}_{j.})}{ (s_W \sqrt{2/n} \;)}
$$
\item We reject $H_0$ when $|t| \geq t_{\alpha/2}$,
 This is equivalent to rejecting $H_0$, for a pair of ($i,j$)
whenever
$$
  |\bar{y}_{i.} - \bar{y}_{j.}| \geq t_{\alpha/2} \,
  s_W \, \sqrt{2/n } \ .
$$
\item The right hand member of this inequality is not a function of $i$ or $j$.
It is constant for a specified $\alpha$ and $n$, and is called the {\textcolor{magenta}{Least Significant
Difference or LSD}}.
\end{itemize}
\end{frame}



\begin{frame}
\begin{itemize}
\item  Once the LSD is calculated,  doing the tests for the pairs of differences of the
form $H_0: \mu_i - \mu_j = 0$ is simple:  Form all possible
absolute differences $|\bar{y}_{i.} - \bar{y}_{j.}|$ and reject
the corresponding $H_0$ if this difference exceeds or  equals the LSD.
%************************************************************************
\item Testing the hypotheses is thus easy, but {\textcolor{magenta}{reporting the
results}} of all those tests can be messy. For $t$ means, there
are $t(t-1)/2$ differences to test.
\item To minimize the number of comparisons we need to make, first 
arrange  the $\bar{y}_{i.}$'s  {\textcolor{magenta}{ordered smallest to largest}}
in value.  If we use the notation $\bar{y}_{(i)}$
for the $i$-th smallest $\bar{y}$, the ranked means may be
represented as 
$
  \bar{y}_{(1)} \leq  \bar{y}_{(2)} \leq  \bar{y}_{(3)}
  \leq \cdots \leq  \bar{y}_{(t)}
$
\end{itemize}
\end{frame}



\begin{frame}
\begin{itemize}
\item For e.g., $\bar{y}_{(1)}$ might be $\bar{y}_7$ if $\bar{y}_{7}$ is
the smallest;  Now note that if the  difference  $\bar{y}_{(t)}- \bar{y}_{(1)}$,
for example, does not exceed the LSD, then all the  differences
$\bar{y}_{(t)}- \bar{y}_{(2)}$,  $\bar{y}_{(t)}- \bar{y}_{(3)}\ldots$,
$\bar{y}_{(t)}- \bar{y}_{(t-1)}$ will not exceed the LSD.
\item It follows that in this case we are spared from computing all the 
above differences and comparing them to the LSD. The following 
procedure is based on this idea:
\item First write the
ordered  means on a line identified by their corresponding treatment names
above them.
\item For Example -- we might have\\[.1in]
\hspace*{1in}\begin{tabular}{lllll}
       trt5 & trt3 & trt1 & trt4 & trt2 \\[.1in]
  9.5 & 10.5 & 11.6 & 12.2 & 13.5\\
       \end{tabular}
       
\end{itemize}
\end{frame}


  
\begin{frame}
\begin{itemize}
\item Take each column in turn, and on a
separate line below the list, starting from column 1 connect the
means by underlining those pairs of means that are not significantly 
different from the mean in the current column, in the following way.
\item Start the comparison of the mean $\bar{y}_{(1)}$
with the mean on the last column $\bar{y}_{(t)}$. We know that if this pair is 
{\tt less than the LSD value}, then none of the differences $|\bar{y}_{(1)}-\bar{y}_{(t-1)}|$
will exceed the LSD value. If so, underline the means connecting $\bar{y}_{(1)}$ with $\bar{y}_{(t)}$
\item Otherwise, move left to the next largest mean $\bar{y}_{(t-1)}$ and compare
 $\bar{y}_{(1)}$ with $\bar{y}_{(t-1)}$, and so on.
\end{itemize}
\end{frame}

 
\begin{frame}
\begin{itemize} 
\item Begin underlining at that  column the the difference is found
to be {\tt less than the LSD value} and extend all the way 
to the left to column 1 (or the column where you started) 
\item This line 
implies that those means that are connected with this line are 
not  significantly different from the mean in column 1 and all means between.
\item Now restart at column 2 (i.e., $\bar{y}_{(2)}$ and repeat the procedure 
the same way as above. The new set of underlines will be displayed in a separate line.
 For Example -- we might have\\[.1in]
\hspace*{2in}\begin{tabular}{lllll}
       trt5 & trt3 & trt1 & trt4 & trt2 \\[.1in]
  9.5 & 10.5 & 11.6 & 12.2 & 13.5\\[-0.2in]
       \\ \cline{1-4}\\[.0001in]
        \cline{2-4}\\[.0002in]
        \cline{3-4}\\[.0003in]
        \cline{4-5}
       \end{tabular}
\end{itemize}
\end{frame}



\begin{frame}
\frametitle{Example}
\begin{itemize}
\item Supposed that the computed sample means of six treatments with equal
sample size 5 (i.e. $n=5$) are:\\[-.3in]
$$\bar{y}_{1.} = 505, \ \bar{y}_{2.} = 528, \ \bar{y}_{3.} = 564, \
\bar{y}_{4.} = 498, \ \bar{y}_{5.} = 600, \ \bar{y}_{6.} = 470
$$
\item Since MSE $=s^2_W=2,451$ with 24 d.f. Thus the LSD is: \\[.15in]
\hspace*{1.5in}LSD $=2.064\sqrt{2(2451)/5}=64.63$.
\item Ordered smallest to largest, the means are:\\[.1in]
$
 \bar{y}_{6.},\,\bar{y}_{4.},\,\bar{y}_{1.},\,\bar{y}_{2.},\,\bar{y}_{3.},\,\bar{y}_{5.},\equiv 470,\, 498,\, 505,\, 528,\, 564,\, 600
$
\item Prepare table to be used in the underlining procedure:
\begin{center}
\begin{tabular}{lllllll}
& trt6 & trt4 & trt1 & trt2 & trt3 & trt5 \\
&  470 & 498 & 505 & 528 & 564 & 600\\[.5in]
\end{tabular}
\end{center}
\end{itemize}
\end{frame}


\begin{frame}
\frametitle{Example - cont.}
\begin{itemize}
\item Using LSD $ = 64.63$,
underlining procedure  is done as follows:
\hspace*{1in}\begin{tabular}{llllll}
       trt6 & trt4 & trt1 & trt2 & trt3 & trt5 \\[.1in]
  470 & 498 & 505 & 528 & 564 & 600 \\[-0.2in]
       \\ \cline{1-4}\\[.0001in]
        \cline{2-4}\\[.0002in]
        \cline{3-5}\\[.0003in]
        \cline{4-5}\\[.0004in]
        \cline{5-6}\\
       \end{tabular}
\item Deleting the superfluous lines we have:
\hspace*{1in}\begin{tabular}{llllll}
       trt6 & trt4 & trt1 & trt2 & trt3 & trt5 \\[.1in]
  470 & 498 & 505 & 528 & 564 & 600 \\[-0.2in]
       \\ \cline{1-4}\\[.0001in]
        \cline{3-5}\\[.0002in]
        \cline{5-6}\\[.0003in]
       \end{tabular}
\end{itemize}
\end{frame}


\begin{frame}
\frametitle{Example - cont.}
These may lead to one or more of the following conclusions:\\[.1in]
\begin{itemize}
\item $\mu_6, \mu_4, \mu_1, \mu_2$ are significantly different from
         $\mu_5$.
\item $\mu_6, \mu_4$ are significantly different from $\mu_3$.
\item  None of $\mu_6, \mu_4, \mu_1$ is significantly different from $\mu_2$.
\item None of $\mu_6, \mu_4$ is significantly different from
           $\mu_1$.
\item $\mu_6$ is not significantly different from $\mu_4$.
\end{itemize}
\end{frame}





\begin{frame}
 When sample sizes are not equal the above procedure is not feasible.
In this case, we may construct confidence intervals for all pairs of
differences $\mu_i - \mu_j $ using
$$
 \bar{y}_{i.} - \bar{y}_{j.} \pm  t_{\alpha/2} \, s_W \, \sqrt{\frac{1}{n_i} + \frac{1}{n_j}}
$$
where $  t_{\alpha/2}$ is again the percentile from the 
t-table with degrees of freedom same as that of the within mean square
$s_W^2$.
\end{frame}


\begin{frame}
\frametitle{Important comments regarding Multiple Comparison procedures}
\begin{itemize}
\item The protected LSD has a {\textcolor{magenta}{per-comparison error rate}} 
of $\alpha$, i.e., the probability of a Type I error is $\alpha$ 
for any single comparison (or test).  However, as we already
discussed, the overall error rate when multiple tests are made 
can be much larger than $\alpha$, i.e.,  the probability of making 
one or more Type I errors exceeds $\alpha$.
\item The {\textcolor{magenta}{protected}} part involves making sure that 
$H_0: \mu_1 = \mu_2 = \cdots = \mu_t$ is tested using the 
analysis of variance F-test prior to using the multiple comparison
procedure. 
\end{itemize}
\end{frame}



\begin{frame}
\frametitle{Important comments regarding Multiple Comparison procedures - cont.}
\begin{itemize}
\item The LSD analysis is carried out only when $H_0$ is
rejected.  There is some evidence, based on 
simulation studies that the {\textcolor{magenta}{experimentwise error rate}} 
for protected LSD may be near $\alpha$.  
\item {\textcolor{magenta}{The experimentwise error
rate}} is the probability of observing an experiment with one or more 
pairwise comparisons falsely declared significant.
\item Protected LSD is not a very conservative method.  We would not be
surprised to see it falsely declare several pairwise comparisons
significant in an experiment involving several treatments when all
possible differences are tested.
\end{itemize}
\end{frame}


\begin{frame}
\frametitle{Important comments regarding Multiple Comparison procedures - cont.}
\begin{itemize}
\item This procedure should not be used to make tests suggested after
the experiment has been conducted and the sample means computed.
At the planning stage of an experiment, the experimenter must
state all questions that needs to  be answered in terms of
possible comparisons. These comparisons are called 
{\textcolor{magenta}{pre-planned or apriori comparisons}}. 
\item Instead of pre-planned
comparisons, a part of the plan for the experiment may require 
testing all differences or only some of them.  The intent of LSD is to 
not to perform all  paiwise comparisons routinely.
\end{itemize}
\end{frame}



\begin{frame}
\frametitle{Important comments regarding Multiple Comparison procedures - cont.}
\begin{itemize}
\item In any case, it is not recommended that any kind of comparisons 
be devised after first looking at the $\bar{y}$'s.
The problem with testing based on comparisons suggested by looking
at the data is that it changes the $\alpha$ level of the test i.e,
Type I error rate is not controlled at the specified  $\alpha$
level anymore.
\item For example --- an extreme case --- say you look at the sample
means and see that the largest is much greater than the smallest,
so you decide to test their difference for significance.  On
average --- across experiments --- you will seldom fail to
reject $H_0$ when you do this, so the Type I error rate is
probably not $\alpha$.
\end{itemize}
\end{frame}


\begin{frame}
\frametitle{Tukey's W Procedure} 
\begin{itemize}
\item This method for comparing all possible pairs is {\textcolor{magenta}{more conservative
than LSD}} (i.e., it tends to be more resistant to falsely declaring
significance.)
\item The method is based on comparisons of
$|\bar{y}_{i.} - \bar{y}_{j.}|$ to the value
$$
  W = q_\alpha(t,\nu) \sqrt{\frac{s_W^2}{n}}
$$
where $s_W^2$ is the mean square within samples all of size $n$,
$\nu$  is the degrees of freedom for $s_W^2$, $t$ is number of
population means, $\mu_i$ compared, and $\alpha$ is the chosen 
significance level.
\end{itemize}
\end{frame}




\begin{frame}
\frametitle{Tukey's W Procedure - cont.} 
\begin{itemize}
\item Then if
$$
  |\bar{y}_{i.} - \bar{y}_{j.}| \geq W 
$$
we declare that the mean pair $\mu_i$ and $\mu_j$ are
significantly different.  
\item The value of $q_\alpha(t,\nu)$ is found in 
Table 10 (Handout).  The table gives $q_\alpha(t,\nu)$ for
either $\alpha=0.05$ or $\alpha=0.01$. 
\item  The the sample means are ordered smallest to largest as before. 
Then make all possible pairwise comparisons using
the value of $W$ and underlining method may then be used to display results. 
\end{itemize}
\end{frame}



\begin{frame}
\frametitle{Example} 
\begin{itemize}
\item The ANOVA table resulting from an experiment involving 6 treatments and $n=5$
per treatment is:\\[.1in]
\begin{tabular}{lrrrr}
Source of Variation & DF &
\multicolumn{1}{c}{SS} &
\multicolumn{1}{c}{MS} &
\multicolumn{1}{c}{F} \\ \hline
Between Treatments & 5 & 847.05 & 169.41 & 14.37 \\
Within Treatments & 24 & 282.93 & 11.79 \\ \hline
Total & 29 & 1129.98
\end{tabular}
\vspace*{.05in}

\item The sample treatment means are:
$\bar{y}_{1.}= 28.8, \bar{y}_{2.} = 24.0, \bar{y}_{3.} = 14.6, \bar{y}_{4.} = 19.9, \bar{y}_{5.} = 13.3, \bar{y}_{6.} =18.7 $\\[.1in]
\item The ordered means table: \\[-.5in]
\begin{center}
\begin{tabular}{lllllll}
& trt5 & trt3 & trt6 & trt4 & trt2 & trt1 \\
& 13.3 & 14.6 & 18.7 &  19.9 & 24.0 & 28.8\\
\end{tabular}
\end{center}
\end{itemize}
\end{frame}

\begin{frame}
\frametitle{Example - cont.} 
\begin{itemize}
\item From Table 10, $q_{.05}(6,24) = 4.37$, so 
$$W = 4.37 \sqrt{\frac{11.79}{5}} = 6.7$$
\item The underlining procedure gives:\\[.01in]
\hspace*{1in}\begin{tabular}{llllll}
       trt5 & trt3 & trt6 & trt4 & trt2 & trt1 \\[.1in]
  13.3 & 14.6 & 18.7 & 19.9 & 24.0 & 28.8 \\[-0.2in]
       \\ \cline{1-4}\\[.0001in]
        \cline{2-4}\\[.0002in]
        \cline{3-5}\\[.0003in]
        \cline{4-5}\\[.0004in]
        \cline{5-6}\\
       \end{tabular}
\end{itemize}
\end{frame}

\begin{frame}
\frametitle{Example - cont.} 
\begin{itemize}
\item  Means that have an underline in common  are declared 
not significantly different from each other.
\item Thus we find that $\mu_5,\, \mu_3,\, \mu_6,$ and $ \mu_4 $ are not
different from $\mu_1$, and $\mu_5$ and $ \mu_3 $ are not different 
from $\mu_2$.
\item Just as with LSD, Tukey's method can be used when sample sizes are
not all the same, but the above procedure is not feasible.  In this 
case we may construct confidence intervals for all pairs of 
comparisons $\mu_i - \mu_j $.  Its form is
$$
\bar{y}_{i.} - \bar{y}_{j.} \pm q_\alpha (t,v) \sqrt{\frac{s_W^2}{2}
\left(\frac{1}{n_i} + \frac{1}{n_j} \right)        }
$$
The value of $q_\alpha (t,\nu)$ from Table 10 is obviously the same for 
all comparisons as well as $s_W^2$. \\[.1in]
\end{itemize}
\end{frame}




\begin{frame}
\frametitle{Scheffe's Procedure}
\begin{itemize}
\item This procedure is ultra conservative.  It controls experimentwise 
error rate.  The probability of observing an experiment 
with one or more contrasts (from the set of all possible 
contrasts) falsely declared significant is the selected
$\alpha$. 
\item Scheffe's method can be used to test all possible differences of
means (recall that simple differences are contrasts).  However it
is usually used where contrasts that are not all
simple differences are to be tested together with any pairwise
differences.
\item  To test $
  H_0: \ell = \sum_i a_i\mu_i = 0 \ \rm{vs.}~ H_a: \ell \neq 0$\\[.1in]
 we base the test statistic on the estimate $\hat{\ell} = \sum a_i \bar{y}_{i.}$
\end{itemize}
\end{frame}



\begin{frame}
\frametitle{Scheffe's Procedure - cont.}
\begin{itemize}
\item Compute the quantity $S$ based on a $F$-distribution as
$$  S = \sqrt{\hat{V}(\hat{\ell})} \ \sqrt{(t-1)F_{\alpha,df_1,df_2}}$$
\item Here $ df_1 = t-1$, and  $df_2 = \nu$ 
\item The variance estimate of $\hat{\ell}$ is  $\hat{V}(\hat{\ell}) = s_W^2 \sum_i \frac{a_i^2}{n_i}$
where $s_W^2$ has $\nu$ degrees of freedom.
\item We reject 
$H_0$ when $|\hat{\ell}| > S$.  
\item Now the underlining procedure is applied in the same way as described for the LSD or
Tukey procedures.
\end{itemize}
\end{frame}
\end{document}
 






